{
\XeTeXlinebreaklocale "th_TH"
\XeTeXlinebreakskip = 0pt plus 1pt
\thaifont
\thaiabstract

การวิจัยครั้งนี้มีวัตถุประสงค์  เพื่อศึกษาความพึงพอใจในการให้บริการงานทั่วไปของสานักวิชา พื้นฐานและภาษา เพื่อเปรียบเทียบระดับความพึงพอใจต่อการให้บริการงาน ทั่วไปของสานักวิชาพื้นฐานและภาษา ของนักศึกษาที่มาใช้บริการสานักวิชาพื้นฐานและภาษา สถาบัน เทคโนโลยีไทย-ญี่ปุ่น จาแนกตามเพศ คณะ และชั้นปีที่ศึกษา เพื่อศึกษาปัญหาและข้อเสนอแนะของ นักศึกษามาเป็นแนวทางในการพัฒนาและปรับปรุงการให้บริการของสานักวิชาพื้นฐานและภาษา
เนื่องจากในยุคสมัยนี้โลกเราได้เข้าสู่ยุคของข้อมูลมหัต (Big data) ซึ่งบริษัทต่าง ๆก็เริ่มขับเคลื่อนธุรกิจด้วยข้อมูลแล้ว  โดยธุรกิจที่มีข้อมูลขนาดใหญ่อย่างที่เห็นง่ายๆเลยก็จะเป็น เว็บไซต์ที่เปิดบริการให้ขายของออนไลน์ ซึ่งในเว็บไซต์นั้นก็จะมีข้อมูลสินค้ามากมายและ สิ่งที่เราได้พบคือรายละเอียดสินค้าหลายๆอย่างไม่มีความสอดคล้องกับรูปภาพสินค้าหรือไม่ชัดเจนพอและ เนื่องด้วยข้อมูลที่มีขนาดใหญ่เกินกว่าที่จะจัดการด้วยแรงงานคนได้และ ยังไม่มีเครื่องมือใด ๆ ในการตรวจจับคำผิดเหล่านี้ได้ เราจึงทำโครงงานนี้ขึ้นมาเพื่อที่จะแก้ไขปัญหาเหล่านี้\par
ทางเราจึงคิดที่จะทำระบบปัญญาประดิษฐ์หรือ AI มาแก้ไขปัญหาในส่วนนี้ โดยทำ Image captioning ที่ใช้ Convolutional neural network และ Long short-term memory ซึ่งทางคณะผู้จัดทำก็จะมีการเปรียบเทียบกับสถาปัตยกรรมของโมเดลอื่น ๆ เพื่อสร้างโมเดลที่มีประสิทธิภาพมากที่สุด ซึ่งถ้าโครงงานนี้ประสบความสำเร็จจะสามารถตรวจสอบคุณภาพของคำอธิบายสินค้าในเว็บไซต์ขายสินค้าออนไลน์ ทำให้คุณภาพและภาพลักษณ์ต่อสังคมของเว็บไซต์นั้นดูดีขึ้นได้ และนอกจากนี้โครงงานนี้ยังสามารถเป็นพื้นฐานให้แก่ผู้ที่สนใจในการทำ image captioning ที่เป็นภาษาไทยและนำไปประยุกต์ใช้กับงานของตนได้อีกด้วย

\begin{flushleft}
\begin{tabular*}{\textwidth}{@{}lp{0.8\textwidth}}
 & \\
\textbf{คำสำคัญ}: & การชุบเคลือบด้วยไฟฟ้า / การชุบเคลือบผิวเหล็ก /  เคลือบผิวรังสี
\end{tabular*}
\end{flushleft}
\endabstract
}
